\documentclass[11pt,a4paper]{article}
\usepackage[margin=2.5cm]{geometry}
\usepackage{amsmath,amssymb}
\usepackage{booktabs}
\usepackage{graphicx}
\usepackage[colorlinks=true,linkcolor=blue,citecolor=blue,urlcolor=blue]{hyperref}
\usepackage{natbib}
\bibliographystyle{plainnat}

\newcommand{\mre}{\mathrm{MRE}}
\newcommand{\kev}{\,\mathrm{keV}}
\newcommand{\kms}{\,\mathrm{km\,s^{-1}}}
\newcommand{\code}[1]{\texttt{#1}}

\title{Neural emulation of SPEX single-element CIE spectra:\\
architecture, training strategy, and evaluation\\
\large Technical report on design decisions, including discarded approaches}
\author{D.~Huppenkothen \and Claude (Anthropic)\thanks{Drafted with
Claude~Fable~5; all results produced by the \code{spexai} codebase,
branch \code{neural\_operator}.}}
\date{2026 July 13}

\begin{document}
\maketitle

\begin{abstract}
We document the design of a neural emulator for velocity-broadened,
single-element collisional ionisation equilibrium (CIE) X-ray spectra
computed with SPEX \citep{kaastra96}, developed on the iron ($Z=26$)
model as testbed. The final configuration -- a coordinate-based
``operator'' network with Fourier-feature energy encoding, FiLM
temperature conditioning, a dedicated line-amplitude head, error-prioritized
example sampling, Polyak weight averaging, and a long training schedule --
reaches a mean relative flux error of $0.13\%$ on held-out SPEX spectra
(99.2\% of test spectra below 1\%), at 21.5\,ms per spectrum. We lay out
each architecture and training decision in the order it was made, with
the quantitative evidence behind it, and we give equal space to the
approaches that were tried and discarded: a Sobolev derivative-matching
loss, PCA+GP emulation, a monolithic $(T,v)$ emulator of broadened
spectra, enlarging the line head, synthetic-data augmentation on dense
grids, and unweighted per-channel instrument metrics. We further
describe validation at the level of real instrument responses
(XRISM/Resolve, Chandra ACIS and HETG), the physics argument for
per-element velocity broadening, an active-learning scheme for placing
new SPEX evaluations, and a feasibility analysis for full 30-element
inference within a 200\,ms budget.
\end{abstract}

\section{Problem setting and scope}
\label{sec:setting}

The SPEX CIE model represents an optically thin plasma as a sum of
per-element spectra: the observed model spectrum is a linear combination
$\sum_i A_i f_i(E; T)$ of single-element spectra $f_i$ weighted by
abundances $A_i$, subsequently velocity-broadened, redshifted, and
folded through an instrument response. Fitting such models with
microcalorimeter-resolution data is limited by the cost of evaluating
$f_i$; we therefore emulate each element's spectrum with a neural
network. This report covers the single-element ($Z=26$, iron) testbed:
data, architecture, training, broadening, and evaluation, in the order
decisions were made. Scaling to all 30 elements uses an automated sweep
(\code{scripts/run\_all\_elements.py}) with the recipe fixed by the
experiments below; the sweep and a failure mode it exposed on low-$Z$
elements are in Section~\ref{sec:multielement}.

\paragraph{Data.} $14{,}246$ SPEX spectra on an adaptive energy grid of
$24{,}212$ bins over $0.1$--$12\kev$, with temperatures
$T = 0.50$--$19.94\kev$; split $81/9/10$ into train/validation/test
($11{,}539$ / $1{,}282$ / $1{,}425$ spectra) with a fixed seed. Targets
are $\log_{10}$ flux densities; bins with $\log_{10} f < -10$ (``the
floor'') are treated as empty. Validation and test sets contain only
original SPEX spectra in every experiment, including all augmentation
studies.

\paragraph{Headline metric.} For spectrum $s$, predicted log-density
$p_{sb}$ and truth $t_{sb}$,
\begin{equation}
\varepsilon_{sb} = \bigl|10^{\,d_{sb}} - 1\bigr|,\qquad
d_{sb} = \operatorname{clip}\bigl(p_{sb} - \max(t_{sb},\,-10),\,-4,\,4\bigr),
\label{eq:eps}
\end{equation}
which equals the relative error in linear flux, and
\begin{equation}
\mre_s = \frac{1}{|V_s|}\sum_{b\in V_s}\varepsilon_{sb},\qquad
V_s = \{\,b : t_{sb} > -10\,\}.
\label{eq:mre}
\end{equation}
We report the mean and median of $\mre_s$ over the test set and
\emph{yields}, the percentage of spectra with $\mre_s$ below $1\%$
(and $0.1\%$). Empty bins are scored separately as \emph{floor
violations} (fraction of empty bins predicted above the floor, and the
worst excess in dex). Line bins ($5.9\%$ of the grid; flux exceeding
twice a running-median continuum) and continuum bins are also scored
separately.

\section{Architecture}
\label{sec:arch}

The emulator is a coordinate-based network in the style of neural
operators \citep{lu21deeponet,li21fno} and follows
\citet{ricketts26}: the spectrum is a continuous function
$\log_{10} f(x;\theta)$ of the normalised log-energy coordinate $x$,
conditioned on normalised parameters $\theta$ (here $\log_{10} T$).
Components, each validated by ablation (Section~\ref{sec:ablation}):

\begin{enumerate}
\item \textbf{Fourier-feature encoding of energy}
      \citep{tancik20}: $x \mapsto [x, \cos 2\pi f_k x, \sin 2\pi f_k x]$
      with $K$ log-spaced frequencies $f_k \in [0.25, f_{\max}]$. A
      coarse-to-fine curriculum weight per channel activates high
      frequencies progressively during the first
      $30\%$ of training, preventing early high-frequency overfitting.
\item \textbf{FiLM conditioning} \citep{perez18}: a small MLP maps
      $\theta$ to per-layer scale/shift $(\gamma,\beta)$ modulating the
      trunk, rather than concatenating $\theta$ to the input.
\item \textbf{Trend head}: a linear-in-$x$ term $a(\theta)\,x +
      b(\theta)$ factoring out the global spectral slope.
\item \textbf{Line head}: the line-bin set is precomputed from the
      training data (union over a temperature-ordered sample; $3{,}980$
      lines for Fe, deterministic given the cache). Each line has a
      learned embedding; a shared MLP maps [embedding, conditioning]
      to a $\log_{10}$ line flux, which is flux-added
      ($\operatorname{logaddexp}$) to the trunk output, so the trunk
      models only continuum and pseudo-continuum. The head stores line
      energies and native bin widths, so integrated line fluxes can be
      re-deposited exactly onto arbitrary instrument grids
      (\code{deposit()}), a property used heavily in
      Sections~\ref{sec:broadening} and \ref{sec:instruments}.
\item \textbf{Per-bin target normalisation}: the network predicts the
      $z$-scored residual of $\log_{10} f$ per training bin; the static
      line forest lives in the normalisation constants and the trunk
      models $\mathcal{O}(1)$ deviations.
\end{enumerate}

The training loss is a Huber loss on $\varepsilon_{sb}$
(Eq.~\ref{eq:eps}) over $P$ randomly sampled grid points per spectrum,
plus a one-sided floor penalty $10^{\operatorname{relu}(p_{sb}+10)}-1$
on empty bins (the model is pushed onto the floor but never penalised
for matching it), plus an early $\log$-space MSE stabiliser ramped to
zero over the first $20\%$ of steps. Optimiser: AdamW
\citep{loshchilov19} with gradient clipping at unit norm and a cosine
schedule with linear warmup \citep{loshchilov17}.

\subsection{Ablation study (what established the components)}
\label{sec:ablation}

An eight-variant ablation at fixed budget (20k steps, $\sim$1M
parameters) established, in order of effect size: Fourier features are
essential (without them, line MRE is 76\%); FiLM is the second-largest
factor (yield@1\% drops from 88.8 to 60.8 without it); a dedicated
line head gives the best line-bin accuracy of the operator family.
A multi-resolution grid encoding \citep[Instant-NGP style;][]{muller22}
performed like the mid-pack variants but evaluated slower, and was not
pursued.

\paragraph{Discarded: Sobolev (derivative-matching) loss.} A loss term
matching finite-difference slopes between sampled points, natural for
smooth reflection spectra, \emph{hurts} here: removing it improved MRE
from $0.91\%$ to $0.52\%$. The CIE spectrum is a forest of
single-bin spikes; penalising slope mismatch on subsampled points
biases the model against exactly the structure it must fit. All
subsequent configurations omit it.

\paragraph{Retained as a deployment option: fixed-grid MLP.} The
original formulation (an MLP from $T$ directly to all $24{,}212$ bins)
is competitive in accuracy ($0.53\%$) and $30\times$ faster
($0.2$\,ms/spectrum) at $15\times$ the parameters. Because CIE line
energies never move, the standard argument against fixed grids does not
apply; this becomes the distillation target for high-throughput
inference (Section~\ref{sec:inference}).

\section{Hyperparameter search}
\label{sec:hpo}

A two-stage random search \citep{bergstra12} over 24 configurations
(learning rate, batch, points per spectrum, trunk width$\times$depth,
activation, $K$ and $f_{\max}$, line-embedding size, trend head,
per-bin normalisation, stabiliser weight, curriculum length): stage one
at 3k steps, top four re-trained at 20k. Winner \code{t04}: GELU
$384\times5$ trunk, $K=512$, $f_{\max}=4000$, learning rate $3\times
10^{-3}$, batch 128, 2048 points/spectrum, line\_dim 16, trend and
bin-norm on (1.65M parameters): $0.31\%$ test MRE, $94.7\%$ yield@1\%.
Validation ranking generalised cleanly to the test set.

Two findings shaped later work. First, \emph{training length dominated
every hyperparameter} ($0.82\%\to0.31\%$ from 3k$\to$20k steps).
Second, the largest configuration (\code{t16}, 3.6M parameters) was
competitive at 3k steps but \emph{diverged} at 20k at the shared
learning rate of $3\times10^{-3}$: the search had coupled model size to
an unstable learning rate, and the apparent conclusion ``larger models
do not help'' was an artefact (revisited in Section~\ref{sec:tier2}).

\section{Classical baselines (and a discarded one)}
\label{sec:baselines}

Per-bin interpolation across training temperatures is the natural
classical competitor with one physical parameter. Monotone cubic
interpolation \citep[PCHIP;][]{fritsch80} and linear interpolation in
$(\log T,\log f)$ achieve $0.003\%$ MRE on the full training grid --
two orders of magnitude better than any neural variant -- and $0.067\%$
with only 300 training spectra. The emulator's justification is
therefore never raw single-element accuracy but speed, differentiability,
memory (4\,MB vs 1.1\,GB), and extension to joint multi-parameter
spaces where dense tabulation is impossible.

\paragraph{Discarded: PCA + Gaussian process emulation.} The
Speculator-style approach \citep{alsing20} -- PCA on $\log f$, GP
regression on component amplitudes -- saturates at its truncation error
($0.2\%$ with 64 components) and degrades catastrophically at small
$n_{\rm train}$ (9--17\% MRE). Rejected as both baseline and generator
for data augmentation; PCHIP was used instead throughout.

\section{Velocity broadening}
\label{sec:broadening}

\subsection{Physics: why per-element broadening is correct}

Turbulent/bulk velocity broadening acts as a convolution with a kernel
$K_v$ in $\ln E$; convolution is linear, so for a shared kernel,
broadening per element and summing is \emph{exactly} equivalent to
broadening the summed spectrum. Moreover the physically correct total
line width is per-ion,
\begin{equation}
\sigma_i^2 = \sigma_{\rm turb}^2 + \frac{k T}{A_i m_p},
\end{equation}
with the thermal term scaling as $1/\sqrt{A_i}$
\citep{rybicki79} -- the decomposition used in microcalorimeter
analyses \citep{hitomi16,hitomi18} -- so a per-element architecture is
\emph{more} physical than broadening the total: each element can carry
its own thermal width at zero cost. Caveats where the convolution
picture itself breaks (multiphase kinematics, resonant scattering) are
orthogonal to the per-element choice.

\subsection{Benchmark of broadening implementations}

Against exact erf-integral broadening of true spectra: FFT convolution
on a uniform $\ln E$ grid is accurate to $0.03$--$0.06\%$ over
$100$--$1000\kms$ at 10--30\,ms (CPU); the \emph{hybrid} pipeline
(emulate unbroadened $\to$ FFT-broaden the trunk $\to$ add analytic
Gaussian lines with exact per-line fluxes) reaches $0.13$--$0.18\%$;
the legacy sparse-matrix implementation degrades from $0.3\%$ to
$2.8\%$ at $1000\kms$ with $20\%$ of line flux misplaced.
\textbf{Recommendation: replace the sparse method with FFT/hybrid.}

\paragraph{Discarded: monolithic $(T,v)$ emulator of broadened spectra
(option~2, v1).} A two-parameter version of the winning architecture,
trained on FFT-broadened targets with the line head removed (rationale:
broadening removes delta spikes), reached only $2.6\%$ validation MRE
(yield@1\% $=0$) and $1.1$--$1.8\%$ end-to-end -- $8\times$ worse than
the hybrid. The post-mortem identified the causes: (i) velocity is a
\emph{known operator}, not a latent parameter -- asking global FiLM
modulation to express a local convolution discards the strongest
available inductive bias; (ii) at $v_{\min}=50\kms$ lines are still
$\sim$2 native bins wide, so removing the line head deleted the
component handling the sharpest structure while keeping trunk
hyperparameters tuned assuming its presence; (iii) uniform point
sampling on the $83{,}728$-bin uniform grid starves line cores of
supervision; (iv) checkpoint selection never saw the hardest velocities
(evaluation at $\{100,250,600\}\kms$ only); (v) zeroing sub-$10^{-9}$
wings in the targets while penalising above-floor predictions injects
label noise at every line wing. A revised version
(\code{train\_broadened2}) addresses all five: a \emph{Gaussian line
head} with analytic width $\sigma_u = v/c$ and $T$-only amplitudes
(broadening conserves line flux; the factorisation is structural),
line-aware point sampling with offsets $\sim\mathcal{N}(0,\sigma(v'))$
at random velocities, validation at $\{50,100,250,600,1000\}\kms$, and
wing masking with a one-sided cap. These fixes more than halved the
error -- v2 reaches $0.46$--$0.77\%$ end-to-end MRE
(Table~\ref{tab:broad2}) versus v1's $1.1$--$1.8\%$, a
$\sim$2.4$\times$ improvement at every velocity, confirming that each
diagnosed cause was real. But v2 remains $4$--$6\times$ worse than the
hybrid ($0.11$--$0.13\%$) and is the slowest method of all (it evaluates
the trunk on the full uniform grid plus the line deposits), so the
revised option~2 is competent but not competitive.

\begin{table}[ht]
\centering
\begin{tabular}{lcccc}
\toprule
method & $100\kms$ & $300\kms$ & $1000\kms$ & note\\
\midrule
FFT (true input) & 0.03\% & 0.03\% & 0.06\% & accuracy ceiling\\
\textbf{hybrid} & \textbf{0.13\%} & \textbf{0.11\%} & \textbf{0.12\%} &
emulate $+$ analytic broaden\\
emulator\_Tv2 (v2) & 0.49\% & 0.46\% & 0.77\% & monolithic, fixed\\
sparse (true input) & 0.45\% & 0.30\% & 2.8\% & legacy production\\
emulator\_Tv (v1) & 1.17\% & 1.06\% & 1.77\% & monolithic, original\\
\bottomrule
\end{tabular}
\caption{End-to-end broadening MRE vs exact erf broadening. FFT and
sparse take the true spectrum as input (isolating the broadening
operation); the three emulator methods are end-to-end. The revised
monolithic $(T,v)$ emulator (v2) more than halves v1's error but stays
$4$--$6\times$ above the hybrid.}
\label{tab:broad2}
\end{table}

The conclusion is therefore structural, not a training-quality
artifact: even a well-executed monolithic $(T,v)$ emulator cannot beat
the hybrid, because the hybrid performs the convolution \emph{exactly}
(FFT plus analytic line deposit) while the monolith must \emph{learn}
the broadened continuum shape. Broadening belongs \emph{outside} the
network; the hybrid on the best unbroadened backbone is the deployment
choice, and option~2 is settled as the losing branch -- retained here
because v2 is the controlled demonstration that the $8\times$ v1 gap was
structural.

\section{Training-data experiments: active learning and augmentation}
\label{sec:adaptive}

Motivated by simulator cost for future elements, we tested dynamically
generating extra training spectra with the best classical interpolator
(PCHIP), with placement informed by model error -- an active-learning /
adaptive-sampling problem \citep{liu18,wu23,crombecq11,nygaard23,
rogers19}. Design: three arms on a 300-point training grid (evenly
subsampled in $\log T$), all validated and tested exclusively on
original SPEX spectra.

\begin{itemize}
\item \textbf{baseline}: uniform sampling of the grid;
\item \textbf{reweight}: examples sampled $\propto$ a running
      (exponentially averaged) per-spectrum training error, mixed 50/50
      with uniform -- prioritized sampling in the spirit of
      \citet{schaul16};
\item \textbf{adaptive}: reweight $+$ gated generation: every 500 steps,
      new temperatures are drawn with density $\propto$ (local error)$^2$
      \citep[RAD;][]{wu23} from grid intervals whose leave-one-out
      interpolation error \citep[cf.\ CV-Voronoi;][]{liu18} is below
      $1\%$; PCHIP spectra generated there enter batches at $25\%$
      with half loss weight (a low-fidelity teacher in the sense of
      \citealt{kennedy00}).
\end{itemize}

\begin{table}[ht]
\centering
\begin{tabular}{lcccc}
\toprule
arm & test MRE & yield@1\% & synthetic & floor viol.\\
\midrule
baseline & 0.45\% & 91.3\% & 0 & 0.28\%\\
reweight & 0.32\% & 95.8\% & 0 & 0.21\%\\
adaptive & 0.30\% & 96.3\% & 1152 & 0.49\%\\
\midrule
\emph{t04\_long, full 11.5k grid} & \emph{0.31\%} & \emph{94.7\%} &
-- & \emph{0.19\%}\\
\bottomrule
\end{tabular}
\caption{Adaptive-data study ($n_{\rm train}=300$). The adaptive arm
matches the full-grid model with $2.6\%$ of the simulator budget.}
\label{tab:adaptive}
\end{table}

Findings (Table~\ref{tab:adaptive}): most of the gain is
\emph{attention, not data} -- prioritized sampling alone closes 85\% of
the gap at zero cost; the synthetic data adds a small further increment
at the price of doubled floor violations (interpolation smears flux
into empty bins); and 300 SPEX evaluations plus augmentation match the
full-grid model -- a $\sim38\times$ simulator-budget reduction, the
practically important result for new elements. The trust gate rejected
23 intervals (clustered at $T=0.53$--$0.76\kev$, leave-one-out errors
2--4\%): an explicit shortlist of where only \emph{real} SPEX runs
help, closing the loop with simulator-in-the-loop acquisition schemes
\citep{nygaard23,rogers19}.

\paragraph{Discarded: augmentation on dense grids.} On the full grid,
interpolation error ($0.003\%$) lies two orders of magnitude below
emulator error ($0.3\%$): every synthetic spectrum is, at the network's
error resolution, a duplicate of its neighbours. This was argued from
the error budget and confirmed by the Tier-1 result below (the model is
capacity-limited, not data-limited: train and validation MRE are
identical). Augmentation is therefore reserved for sparse-grid regimes
($n_{\rm train}\lesssim$ a few hundred here).

\section{Training strategy: Tier 1 (optimisation) and Tier 2 (capacity)}
\label{sec:tier2}

\subsection{Tier 1}
Three changes, no architecture change, full training grid:
(i) $5\times$ longer schedule (100k steps $\approx$ 1100 epochs) with a
floor on the cosine at $2\%$ of peak learning rate; (ii) Polyak/EMA
weight averaging \citep{polyak92,izmailov18} with decay $0.999$ --
validation and checkpoints use the averaged weights, which also
stabilises the noisy yield-based checkpoint selection; (iii) the
prioritized sampling from Section~\ref{sec:adaptive}. Result:
$0.31\%\to\mathbf{0.16\%}$ MRE, yield@1\% $94.7\to99.0$, yield@0.1\%
$0\to43.8$, floor violations halved. Training and validation MRE were
identical at convergence (0.15\%): no overfitting after 1100 epochs,
implying a capacity limit, and the best checkpoint was the final step,
implying the budget still binds.

\paragraph{Caveat found only by instrument-level tests.} Prioritization
traded away accuracy on the faint high-energy tails of cool ($T\lesssim
2\kev$) spectra -- invisible to Eq.~\ref{eq:mre} (those bins are
floor-masked or negligible in the full-band average) but exposed as a
2--5\% regression in the Chandra grating bands
(Section~\ref{sec:instruments}). Mitigated in Tier~2 by lowering the
prioritized fraction from $0.5$ to $0.4$.

\subsection{Tier 2}
Two capacity hypotheses at equal budget, both with the Tier-1 recipe
and a warmup--stable--decay (WSD) schedule \citep[flat learning rate,
linear decay over the final 15\%; extending a run re-pays only the
decay leg;][]{hu24} :

\begin{table}[ht]
\centering
\begin{tabular}{lcccccc}
\toprule
 & params & MRE & line/cont & yield@1\% & yield@0.1\% & ms/spec\\
\midrule
tier1 (t04 arch) & 1.65M & 0.16\% & 0.21/0.15 & 99.0 & 43.8 & 12.4\\
\textbf{big\_trunk} & 3.3M & \textbf{0.13\%} & 0.19/0.12 & \textbf{99.2}
& \textbf{52.6} & 21.5\\
line\_heavy & 1.9M & 0.21\% & 0.23/0.21 & 98.1 & 25.8 & 12.9\\
\bottomrule
\end{tabular}
\caption{Tier-2 capacity experiments (100k steps, full grid).}
\label{tab:tier2}
\end{table}

\code{big\_trunk} ($512\times6$, $K=1024$, learning rate $10^{-3}$)
wins across the board (Table~\ref{tab:tier2}), formally overturning the
HPO's size conclusion: it was always the learning rate. Scaling is
$\sim25\%$ MRE reduction per parameter doubling with both runs still
budget-clipped, so the next cheap move is the same configuration at
200k steps.

\paragraph{Discarded: enlarging the line head (\code{line\_heavy}).}
Wider line embeddings (48), a wider shared MLP (256), and
Fourier-embedded $T$-conditioning of line amplitudes -- targeting the
dominant line-MRE component -- performed \emph{worse than Tier 1}
($0.21\%$ overall; line MRE $0.23\%$ vs $0.21\%$). The line-error
bottleneck is thus not amplitude-head capacity; the likely residual is
the trunk under-resolving inter-line pseudo-continuum. (Minor caveat:
this arm ran at the default learning rate $3\times10^{-3}$; given it
regressed rather than plateaued, we did not spend a further run on the
learning-rate confound.)

\subsection{Multi-element generalisation: floor-dominated spectra}
\label{sec:multielement}
Fixing the recipe on iron and sweeping the low-$Z$ elements ($Z=2$--$6$)
with the automated per-element sizing exposed a failure mode absent from
iron. Two sharply separated tiers emerged (Table~\ref{tab:multielement}):
helium and carbon reach $\le0.04\%$ MRE at $100\%$ yield on both
thresholds, while lithium, beryllium, and boron land at $1.3$--$9.9\%$
MRE with essentially no spectra under $0.1\%$.

The split is explained entirely by \emph{signal sparsity}, not atomic
number. On the shared $24{,}212$-bin grid the fraction of bins ever above
the $-10$ floor is $\sim97\%$ for He/C (broad continua) but under $1.2\%$
for Li/Be/B -- beryllium carries flux in only $\sim6$ bins per spectrum,
the rest at floor. This is not intrinsic difficulty: classical linear
interpolation in $T$ reaches $\sim10^{-6}$ MRE on all five elements (the
trust gate reports $100\%$ of intervals interpolable), so the targets are
trivially learnable and the poor emulator numbers are a training
pathology.

\begin{table}[ht]
\centering
\begin{tabular}{lccccc}
\toprule
 & $Z$ & signal bins/spec & signal frac. & emul.\ MRE & interp.\ MRE\\
\midrule
He & 2 & $23{,}703$ & $97.9\%$   & $3.2\times10^{-4}$ & $8.8\times10^{-7}$\\
C  & 6 & $23{,}444$ & $96.8\%$   & $4.0\times10^{-4}$ & $8.6\times10^{-7}$\\
B  & 5 & $292$      & $1.2\%$    & $9.9\times10^{-2}$ & $1.1\times10^{-6}$\\
Li & 3 & $15$       & $0.06\%$   & $1.3\times10^{-2}$ & $9.1\times10^{-7}$\\
Be & 4 & $6$        & $0.03\%$   & $5.9\times10^{-2}$ & $9.9\times10^{-7}$\\
\bottomrule
\end{tabular}
\caption{Low-$Z$ sweep before the signal-aware fix: emulator accuracy
tracks signal density, while classical interpolation is uniformly
excellent. Signal bins are per-spectrum bins above the floor (mean over
the test set); MRE is the fraction of Eq.~\ref{eq:mre}.}
\label{tab:multielement}
\end{table}

Two compounding causes. (i)~Uniform point sampling: each step draws
\code{points}$=2048$ of the $24{,}212$ bins uniformly, so for beryllium
the expected number of non-floor bins seen per step is
$2048\times0.0003\approx0.5$ -- the model is supervised on the very bins
it is scored on in only half of all steps. (ii)~Even when a line bin is
drawn, the flat per-point mean in the training loss lets the $\sim2000$
trivial floor bins dominate the gradient. Compounding both, the automated
size classifier keyed on line-bin \emph{count} and routed Li/Be to a
``smooth'' preset that switched the line head \emph{off} -- disabling the
line machinery for precisely the spectra that are almost nothing but
lines.

\paragraph{Fix: signal-aware sampling.} We replace edge-only sampling
with a \emph{signal pool} (union of detected line and
recombination-edge bins) and (a)~force its members into every batch, and
(b)~up-weight them per point so they hold a target share
\code{signal\_frac} of the loss regardless of how few they are. This is
one self-tuning knob that decouples a signal bin's loss \emph{influence}
from its sampled \emph{count}, so the same setting holds across
beryllium's $\sim6$ and boron's $\sim292$ signal bins. The size
classifier is corrected to route on signal \emph{density}: a
floor-dominated spectrum with any lines goes to a new \code{sparse\_lines}
preset (line head on, \code{signal\_frac}~$=0.25$ and learning rate
$10^{-3}$, both set by validation below), which line-bin count
alone cannot distinguish from a genuinely smooth continuum (He). The
broad-continuum elements are unaffected -- their signal pool spans the
band, so the weighting is a near-identity.

\paragraph{Validation on beryllium.} Re-running the worst element
($Z=4$, six signal bins per spectrum) with the corrected classifier and
the \code{sparse\_lines} recipe (line head on, \code{signal\_frac}~$=0.25$,
learning rate $10^{-3}$ -- the inherited $3\times10^{-3}$ reproduced the
old bouncing divergence) transformed it. Held-out test MRE
$5.9\%\to\mathbf{0.013\%}$, yield@1\% $0\to\mathbf{100\%}$, yield@0.1\%
$0\to\mathbf{99.2\%}$, with the line-bin MRE collapsing from $4.2\%$ to
$\sim0$ and floor violations to zero. The run was stable to step $86$k
(the original diverged by $28$k, and its apparent $5.9\%$ was a lucky
checkpoint dip before collapse to MRE $\approx1$). This confirms the
failure was sampling starvation and optimiser instability, not
capacity: six signal bins suffice for $99\%$ of spectra under $0.1\%$
once they are actually supervised. The emulator now sits $\sim140\times$
above the interpolation baseline (from $\sim6\times10^{4}\times$), the
same regime as the well-behaved elements. Confirmation that the fix
transfers to lithium ($Z=3$) and boron ($Z=5$), and that the changed
default \code{signal\_frac} does not regress iron, is pending.

\paragraph{Training throughput.} On one A100, per-step wall time scales
with the preset: $\sim35$\,ms/step for \code{smooth} ($0.2$M params),
$\sim100$\,ms/step for \code{sparse\_lines} ($0.6$M), and
$\sim210$\,ms/step for \code{standard} ($1.6$M). With early stopping most
elements converge in $40$--$95$k steps, so a single full-grid element
trains end to end in roughly $0.7$--$4$ GPU-hours (smooth continua at the
low end, line-rich \code{standard} elements at the high end).

\paragraph{Early-stopping fragility on noisy runs.} Per-element step
count is governed by a smoothed validation-MRE plateau detector, which is
sensitive to the large eval-to-eval noise of the high-learning-rate
\code{standard} runs (val MRE swings of $\pm100$--$300\%$ between evals).
Two failure tails appear. On sodium and aluminium a single upward spike
(Na hit MRE $0.14$ at step $40$k) made the smoothed signal read as
plateaued and triggered the stop at $42$--$48$k; the models then kept
improving $30$--$40\%$ per eval through the forced decay and finished
$\sim60$k still descending, leaving them $3$--$5\times$ worse than their
neighbours ($0.30$--$0.37\%$ MRE, $0\%$ yield@0.1\%, and for Al a
$0.62\%$ floor-violation rate -- all hallmarks of an unconverged fit, the
line and continuum shapes being correct). Conversely fluorine's noise
\emph{prevented} the trigger, running it to the $100$k cap. Both tails are
mitigated by widening the smoothing window (\code{mre\_smooth} $3\to5$) and
the patience (\code{early\_stop\_patience} $5\to8$) so no single eval can
move the decision; the affected elements are being re-run under the new
defaults. Checkpoint-on-lowest-val-MRE ensures the shipped model is the
best seen within a run, so noise never degrades it below the trajectory
-- but a premature stop still forfeits the accuracy the remaining budget
would have bought, which is exactly what cost Na and Al.

\section{Instrument-level validation}
\label{sec:instruments}

Native-grid metrics cannot answer whether the emulator is accurate
\emph{at the resolution an instrument sees}. We evaluate on
(i) resolution-matched grids (bin width $=$ FWHM/2 from published
resolutions) and (ii) the actual detector channel grids, folding both
truth and prediction through the real redistribution matrices: XRISM
Resolve \code{rsl\_Hp\_L\_2025.rmf} \citep[$60000^2$, 93M non-zeros;][]
{ishisaki25}, Chandra ACIS \code{aciss\_aimpt\_cy28.rmf}
\citep{garmire03}, and HETG HEG/MEG grating RMFs \citep{canizares05},
parsed from the OGIP format \citep{george07}. The instrument metric is
Eq.~\ref{eq:mre} applied to integrated flux per channel with validity
mask $F^{\rm true}>0$, band-restricted (Resolve: $1.8$--$12\kev$, the
closed-gate-valve limit). ARF weighting is deliberately excluded
(effective area does not mix channels).

Findings: (i) \emph{median} channel errors (0.03--0.08\%) are better
than native MRE -- rebinning averages independent per-bin errors;
(ii) \code{big\_trunk} is the best model on ACIS and both gratings and
repairs the Tier-1 grating regression; (iii) at Resolve's 0.5\,eV
channels, sub-bin line placement -- exact-bin deposit vs.\ flux split
across a channel edge -- is a real error mode invisible on the native
grid; (iv) unweighted per-channel \emph{means} are ill-conditioned:
rebinning gives every channel a scrap of flux, so near-empty channels
(unbounded relative error on $\sim$zero truth) and the 4/16 coolest
test spectra (negligible in-band flux) dominate, with means swinging by
$30\times$ between models whose medians agree.

\paragraph{Discarded: unweighted per-channel means as a headline
number.} For publication, the flux-weighted error
\begin{equation}
\mre^{\rm flux}_s
= \frac{\sum_c F^{\rm true}_{sc}\,\varepsilon_{sc}}{\sum_c F^{\rm true}_{sc}}
= \frac{\sum_c |F^{\rm pred}_{sc}-F^{\rm true}_{sc}|}{\sum_c F^{\rm true}_{sc}}
\end{equation}
(the misplaced-flux fraction) is immune to empty-channel pathology by
construction; implementation is a planned change. Restricting the
$1.8\kev$ Resolve cutoff changed results by $<2\%$ relative -- the
pathology lives inside the science band, not below it.

\section{Inference feasibility for the 30-element model}
\label{sec:inference}

Requirement: forward passes through 30 element models, abundance
combination, broadening, and response folding in $<200$\,ms. Measured
ingredients: $8.3$\,ms/spectrum amortised (batch 256, A100, full native
grid) for the operator model; Resolve RMF folding is memory-bandwidth
bound (93M non-zeros $\approx0.75$\,GB of traffic): 37\,ms measured on
a laptop CPU, $\sim$0.5--1\,ms projected on A100 HBM -- reconciling the
historical observation that the RMF fold dominated CPU pipelines with
its negligible cost GPU-resident. The naive 30-element loop
($\sim$250\,ms compute-bound) is brought within budget by (i)
evaluating the smooth trunk only at instrument-required resolution
(lines deposited analytically; 3--20$\times$), (ii) bf16 and grouped
batched matmuls over the 30 weight sets (2--4$\times$), (iii) folding
once after summation (linearity), with per-element thermal widths as a
single batched FFT. Projected: 30--80\,ms per likelihood
evaluation. For sampling workloads, distillation into fixed-grid heads
(0.2\,ms/spectrum; Section~\ref{sec:ablation}) using the operator
models as teachers reduces the full stack to $\sim$10\,ms; the RMF can
additionally be truncated at $10^{-5}$ of row peak (54\% of non-zeros,
0.007\% flux loss).

\section{Related work (2023--2026) and candidate improvements}
\label{sec:related}

A survey of recent ML-venue literature identifies candidates in three
categories: improvements on deployed strategies, competing
alternatives, and additions for performance per unit compute.

\paragraph{Improvements on deployed strategies.}
(i) Our error-prioritized sampling is a \emph{biased} gradient
estimator -- the mechanism behind the Tier-1 grating regression.
InfoBatch \citep{qin24} performs soft loss-based pruning \emph{with
gradient rescaling of survivors}, provably unbiased, saving 25--40\%
of training cost; it is a small change to our sampler that removes the
bias and the \code{pr\_mix} heuristic simultaneously.
(ii) Maximum Update Parametrization \citep{yang21mup} transfers
learning rates across widths, formally solving the size--learning-rate
coupling that misled our HPO (Section~\ref{sec:hpo}); it is used for
spectral emulators by \citet{rozanski25scaling}, who also report
compute-optimal scaling exponents for this problem class.
(iii) Schedule-Free AdamW \citep{defazio24} replaces our hand-assembled
WSD-plus-EMA combination (no pre-committed budget; averaged-iterate
inference built in), and Muon \citep{jordan24} applies orthogonalized
updates to weight matrices; recent MLP-regression benchmarks find
Schedule-Free, Muon, and AdamW+EMA (our recipe) the three consistent
winners, making the other two cheap drop-in trials.
(iv) FINER's variable-periodic activations \citep{liu24finer} and
Fourier-reparameterized training \citep{shi24} address spectral bias
adaptively, targeting the fixed-$f_{\max}$ ceiling of our encoding;
EVOS \citep{zhang24evos} formalises coordinate-point selection during
INR training (cf.\ our line-aware sampling).

\paragraph{Competing alternatives.}
TransformerPayne \citep{rozanski25} demonstrates, for stellar spectral
emulation specifically, that attention over the wavelength axis breaks
the $\sim$1\% plateau of MLP emulators with better accuracy \emph{and}
data efficiency, by capturing long-range correlations (co-varying lines
of the same ion) that our independent per-line embeddings ignore -- the
best-motivated architecture competitor, modulo heavier inference.
\citet{yang25grids} argue that for dense, regularly sampled signals,
explicit grids plus interpolation beat implicit networks per unit
compute -- consistent with our interpolation baselines -- which
legitimises a grid/low-rank component plus small neural residual as the
Tier-3 architecture, and strengthens the fixed-grid distillation route.
Kolmogorov--Arnold networks \citep{liu24kan} report favourable scaling
on smooth scientific fitting problems, but wall-clock efficiency claims
replicate poorly and nothing suggests an advantage on line-forest
targets (low priority). Attention/state-space operators
\citep[e.g.][]{wu24transolver} become relevant at the joint
multi-element stage, not for 1-D single-element emulation.

\paragraph{Proposed tests.}
These candidates are \emph{not} suited to a single factorial ablation:
Schedule-Free \emph{replaces} WSD+EMA (mutually exclusive alternatives,
not additive switches), muP is testable only as a width sweep, and
InfoBatch interacts with the sampler. The plan is instead:
(1) A/B with 2--3 seeds at t04 scale for InfoBatch-style rescaling
(success criterion: grating regression absent at full prioritization
strength, equal MRE at lower cost);
(2) a three-arm optimizer bake-off (WSD+EMA / Schedule-Free / Muon) at
short budget, winner confirmed once at 100k steps;
(3) a three-width muP sweep with the learning rate tuned only at the
smallest width;
(4) a FINER $\times$ curriculum $2\times2$ (FINER may subsume the
coarse-to-fine schedule). Winners are composed greedily and confirmed
in a single big\_trunk-scale run, with the test set touched only at
that final step.

\section{Summary of discarded approaches}
\label{sec:discarded}

\begin{table}[ht]
\centering
\small
\begin{tabular}{p{4.6cm}p{4.6cm}p{5.2cm}}
\toprule
Approach & Evidence against & Replacement\\
\midrule
Sobolev derivative loss & MRE $0.91\%\to0.52\%$ on removal & plain
relative-error Huber\\
PCA(64)+GP baseline/generator & saturates at 0.2\%; 9--17\% at small
$n$ & per-bin PCHIP\\
Sparse-matrix broadening & 2.8\% MRE, 20\% flux misplaced at
$1000\kms$ & FFT / hybrid\\
Monolithic $(T,v)$ emulator (v1, v2) & v1 $8\times$, v2 (fixed)
$4$--$6\times$ worse than hybrid & hybrid (broaden outside the net)\\
Line-head enlargement & 0.21\% vs 0.16\% (worse incl.\ line MRE) &
trunk scaling (big\_trunk)\\
Uniform batch sampling & 0.45\% vs 0.32\% (reweight) at $n=300$ &
error-prioritized sampling\\
Augmentation on dense grids & interp.\ error $100\times$ below model
error; train$=$val MRE & augmentation only for sparse grids\\
HPO's ``big models diverge'' & t16 stable and winning at
lr $10^{-3}$ & decouple size from learning rate\\
20k-step cosine budgets & best checkpoint $=$ final step, twice &
100k+ steps; WSD schedule\\
Unweighted per-channel means & $30\times$ swings between models with
equal medians & flux-weighted MRE (planned)\\
$(T,v)$-emulator line-head removal & core errors $\pm$2--5\% at all
$v$ & analytic Gaussian line head\\
\bottomrule
\end{tabular}
\caption{Approaches tried and discarded, with the deciding evidence.}
\label{tab:discarded}
\end{table}

\section*{Reproducibility}

Code: \code{spexai} branch \code{neural\_operator}. Trainers:
\code{spexai/train/train\_operator.py} (variants/ablation),
\code{train\_adaptive.py} (prioritized sampling, gated augmentation,
Tier 1/2 recipes), \code{train\_broadened2.py} (revised $(T,v)$
emulator). Benchmarks: \code{scripts/benchmark\_operator.py} (native
grid), \code{benchmark\_broadening.py}, \code{benchmark\_instruments.py}
(resolution-matched $+$ RMF-folded), \code{baseline\_interpolation.py}.
Sweep driver: \code{scripts/run\_all\_elements.py}. The recipe-screening
program of Section~\ref{sec:related} is implemented in
\code{scripts/run\_recipe\_program.py} (14 arms; new options
\code{--optimizer adamw|schedulefree|muon}, \code{--mup},
\code{--pr\_correct}, \code{--activation finer} in \code{train\_adaptive.py}
and \code{spexai/train/optim.py}); confirmation on fresh off-grid PCHIP
probes rather than the test set is \code{scripts/benchmark\_offgrid.py}. Every training run
writes train-vs-validation curves and three-temperature
spectrum/residual/line-zoom figures automatically
(\code{spexai/train/diagnostics.py}). The winning configuration
(\code{big\_trunk}): GELU $512\times6$ trunk, $K=1024$ Fourier
frequencies ($f_{\max}=4000$), FiLM, trend head, per-bin normalisation,
line head (dim 16, hidden 128), batch 128, 2048 points/spectrum, AdamW
at $10^{-3}$ with WSD schedule, EMA decay 0.999, prioritized sampling
with mix 0.4, 100k steps.

\begin{thebibliography}{99}

\bibitem[Alsing et al.(2020)]{alsing20}
Alsing, J., et al.\ 2020, ApJS, 249, 5 (Speculator)

\bibitem[Bergstra \& Bengio(2012)]{bergstra12}
Bergstra, J., \& Bengio, Y.\ 2012, JMLR, 13, 281

\bibitem[Canizares et al.(2005)]{canizares05}
Canizares, C.~R., et al.\ 2005, PASP, 117, 1144 (Chandra HETG)

\bibitem[Crombecq et al.(2011)]{crombecq11}
Crombecq, K., et al.\ 2011, SIAM J.~Sci.~Comput., 33, 1948
(LOLA-Voronoi)

\bibitem[Fritsch \& Carlson(1980)]{fritsch80}
Fritsch, F.~N., \& Carlson, R.~E.\ 1980, SIAM J.~Numer.~Anal., 17, 238

\bibitem[Garmire et al.(2003)]{garmire03}
Garmire, G.~P., et al.\ 2003, Proc.~SPIE, 4851, 28 (Chandra ACIS)

\bibitem[George et al.(2007)]{george07}
George, I.~M., et al.\ 2007, OGIP Memo CAL/GEN/92-002, ``The
Calibration Requirements for Spectral Analysis''

\bibitem[Hitomi Collaboration(2016)]{hitomi16}
Hitomi Collaboration 2016, Nature, 535, 117

\bibitem[Hitomi Collaboration(2018)]{hitomi18}
Hitomi Collaboration 2018, PASJ, 70, 9

\bibitem[Hu et al.(2024)]{hu24}
Hu, S., et al.\ 2024, arXiv:2404.06395 (MiniCPM; WSD schedule)

\bibitem[Ishisaki et al.(2025)]{ishisaki25}
Ishisaki, Y., et al.\ 2025, J.~Astron.~Telesc.~Instrum.~Syst., 11,
042023 (XRISM Resolve)

\bibitem[Izmailov et al.(2018)]{izmailov18}
Izmailov, P., et al.\ 2018, UAI (stochastic weight averaging)

\bibitem[Kaastra et al.(1996)]{kaastra96}
Kaastra, J.~S., Mewe, R., \& Nieuwenhuijzen, H.\ 1996, in UV and X-ray
Spectroscopy of Astrophysical and Laboratory Plasmas, 411 (SPEX)

\bibitem[Kennedy \& O'Hagan(2000)]{kennedy00}
Kennedy, M.~C., \& O'Hagan, A.\ 2000, Biometrika, 87, 1

\bibitem[Li et al.(2021)]{li21fno}
Li, Z., et al.\ 2021, ICLR (Fourier neural operator)

\bibitem[Liu et al.(2018)]{liu18}
Liu, H., Ong, Y.-S., \& Cai, J.\ 2018, Struct.~Multidisc.~Optim., 57,
393 (adaptive-sampling survey)

\bibitem[Loshchilov \& Hutter(2017)]{loshchilov17}
Loshchilov, I., \& Hutter, F.\ 2017, ICLR (SGDR)

\bibitem[Loshchilov \& Hutter(2019)]{loshchilov19}
Loshchilov, I., \& Hutter, F.\ 2019, ICLR (AdamW)

\bibitem[Lu et al.(2021)]{lu21deeponet}
Lu, L., Jin, P., \& Karniadakis, G.~E.\ 2021, Nature Mach.~Intell., 3,
218 (DeepONet)

\bibitem[M{\"u}ller et al.(2022)]{muller22}
M{\"u}ller, T., et al.\ 2022, ACM Trans.~Graph., 41, 102 (Instant-NGP)

\bibitem[Nygaard et al.(2023)]{nygaard23}
Nygaard, A., et al.\ 2023, JCAP, 05, 025 (CONNECT)

\bibitem[Perez et al.(2018)]{perez18}
Perez, E., et al.\ 2018, AAAI (FiLM)

\bibitem[Polyak \& Juditsky(1992)]{polyak92}
Polyak, B.~T., \& Juditsky, A.~B.\ 1992, SIAM J.~Control Optim., 30,
838

\bibitem[Ricketts et al.(2026)]{ricketts26}
Ricketts, B., Hadzi Veljkovic, D., et al.\ 2026, in preparation
% TODO: complete reference when published

\bibitem[Rogers \& Peiris(2019)]{rogers19}
Rogers, K.~K., \& Peiris, H.~V.\ 2019, JCAP, 02, 031

\bibitem[Qin et al.(2024)]{qin24}
Qin, Z., et al.\ 2024, ICLR (oral), arXiv:2303.04947 (InfoBatch)

\bibitem[Defazio et al.(2024)]{defazio24}
Defazio, A., et al.\ 2024, NeurIPS, arXiv:2405.15682
(Schedule-Free learning)

\bibitem[Jordan et al.(2024)]{jordan24}
Jordan, K., et al.\ 2024, ``Muon: an optimizer for hidden layers in
neural networks'', \url{https://kellerjordan.github.io/posts/muon/}

\bibitem[Yang \& Hu(2021)]{yang21mup}
Yang, G., \& Hu, E.~J.\ 2021, ICML; see also Yang et al.\ 2022,
arXiv:2203.03466 (Maximum Update Parametrization / $\mu$Transfer)

\bibitem[Liu et al.(2024a)]{liu24finer}
Liu, Z., et al.\ 2024, CVPR (FINER: variable-periodic activations)

\bibitem[Shi et al.(2024)]{shi24}
Shi, K., et al.\ 2024, CVPR, arXiv:2401.07402
(Fourier-reparameterized training)

\bibitem[Zhang et al.(2024)]{zhang24evos}
Zhang, W., et al.\ 2024, arXiv:2412.10153 (EVOS: evolutionary
coordinate selection for INR training)

\bibitem[Liu et al.(2024b)]{liu24kan}
Liu, Z., et al.\ 2024, arXiv:2404.19756 (Kolmogorov--Arnold networks)

\bibitem[R{\'o}{\.z}a{\'n}ski et al.(2025)]{rozanski25}
R{\'o}{\.z}a{\'n}ski, T., Ting, Y.-S., et al.\ 2025, ApJ,
arXiv:2407.05751 (TransformerPayne)

\bibitem[R{\'o}{\.z}a{\'n}ski \& Ting(2025)]{rozanski25scaling}
R{\'o}{\.z}a{\'n}ski, T., \& Ting, Y.-S.\ 2025, Open J.~Astrophys.,
arXiv:2503.18617 (scaling laws for spectral emulation)

\bibitem[Yang et al.(2025)]{yang25grids}
2025, arXiv:2506.11139 (``Grids often outperform implicit neural
representations'')
% TODO: verify author list

\bibitem[Wu et al.(2024)]{wu24transolver}
Wu, H., et al.\ 2024, ICML (Transolver)

\bibitem[Rybicki \& Lightman(1979)]{rybicki79}
Rybicki, G.~B., \& Lightman, A.~P.\ 1979, Radiative Processes in
Astrophysics (Wiley)

\bibitem[Schaul et al.(2016)]{schaul16}
Schaul, T., et al.\ 2016, ICLR (prioritized experience replay)

\bibitem[Tancik et al.(2020)]{tancik20}
Tancik, M., et al.\ 2020, NeurIPS (Fourier features)

\bibitem[Wu et al.(2023)]{wu23}
Wu, C., et al.\ 2023, Comput.~Methods Appl.~Mech.~Eng., 403, 115671
(RAD/RAR-D adaptive sampling)

\end{thebibliography}

\end{document}
